\documentclass[11pt]{article}
\usepackage[margin=1in]{geometry}
\usepackage{listings}
\usepackage{xcolor}
\usepackage{graphicx}
\usepackage{float}
\usepackage[hidelinks]{hyperref}

\lstset{
    language=C++,
    basicstyle=\ttfamily\small,
    keywordstyle=\color{blue},
    commentstyle=\color{gray},
    stringstyle=\color{red},
    breaklines=true,
    showstringspaces=false,
    tabsize=4,
    numbers=left,
    numberstyle=\tiny,
    frame=tb
}


\title{Laborator 7 - SO}
\author{Rus Lucian-Florin}
\date{May 6, 2026}

\begin{document}

\maketitle

\newpage
\section*{Comenzi laborator}

\begin{figure}[H]
    \centering
    \includegraphics[width=0.8\textwidth]{1.png}
    \caption{Set comenzi laborator 1}
    \label{fig:1}
\end{figure}

\newpage
\section*{Teme}
\subsection*{Problema 2}
\begin{lstlisting}
#include <pthread.h>
#include <semaphore.h>
#include <stdio.h>
#include <stdlib.h>

// structura care ne permite sa transmitem o variabila de prioritate alaturi de mesaj
// pornirea thread-urile intr-o ordine data necesita o schema de sincronizare specifica
// in cazul de fata optam pentru o variabila care sa ne permita acest lucru
typedef struct {
    char msg[24]; // alocare static pentru ca e mai usor
    int  prio;
} context_t;

sem_t thread2_start;

void print_message_function(void *ptr) {
    // daca avem prio 2, asteptam ca semaforul sa fie setat
    if (((context_t *)ptr)->prio == 2) {
        sem_wait(&thread2_start);
    }

    printf("%s ", ((context_t *)ptr)->msg);

    // daca avem prio 1, setam semaforul
    if (((context_t *)ptr)->prio == 1) {
        sem_post(&thread2_start);
    }
}

int main(int argc, char *argv[]) {
    pthread_t thread1, thread2;

    // numar mai mare -> prio mai mica
    context_t msg1 = {
        .prio = 1,
        .msg  = "Hello",
    };
    context_t msg2 = {
        .prio = 2,
        .msg  = "world",
    };

    // initializam cu 0 pentru ca thread2 sa nu porneasca
    sem_init(&thread2_start, 0, 0);

    // trimitem structurile prin cast la void
    pthread_create(&thread1, NULL, (void *)&print_message_function, (void *)&msg1);
    pthread_create(&thread2, NULL, (void *)&print_message_function, (void *)&msg2);

    pthread_join(thread1, NULL);
    pthread_join(thread2, NULL);
    exit(0);
}

\end{lstlisting}

\bigskip
\bigskip
\bigskip
\bigskip
\bigskip
\bigskip

\begin{figure}[H]
    \centering
    \includegraphics[width=0.8\textwidth]{2.png}
    \caption{Output \textbf{hello} inainte si dupa modificare}
    \label{fig:2}
\end{figure}

Putem observa cum, inainte de modificari, comportamentul nu este deterministic, fapt care duce la nevoia de a implementare o schema de sincronizare specifica. Pentru implementare, s-a folosit o structura cu un identificator de prioritate, reprezentat de un numar intreg.

\newpage
\subsection*{Problema 3}
\begin{lstlisting}
#include <pthread.h>
#include <stdio.h>
#include <stdlib.h>

#define ITEMS 10

long buffer[ITEMS];
int  head = 0, tail = 0;

pthread_mutex_t mutex;
pthread_cond_t  thread_condition;
struct timespec delay;

long produce_item(void) {
    long item = random() % 256;
    printf("producing %d\n", item);

    return item;
}

void consume_item(long item) {
    printf("consuming %d\n", item);
}

void producer_function(void) {
    while (1) {
        pthread_mutex_lock(&mutex);
        if ((tail + 1) % ITEMS != head) {
            buffer[tail] = produce_item();
            tail         = (tail + 1) % ITEMS;
        }
        pthread_mutex_unlock(&mutex);

        // facem broadcast pentru ca, desi avem doar un consumer, codul e mai fool proof
        pthread_cond_broadcast(&thread_condition);
        nanosleep(&delay, NULL);
    }
}

void consumer_function(void) {
    while (1) {
        // pthread_mutex_lock(&mutex);
        if (head != tail) {
            consume_item(buffer[head]);
            head = (head + 1) % ITEMS;
        }
        // pthread_mutex_unlock(&mutex);

        // putem sa comentam lock/unlock pentru ca `cond_wait` o sa faca unlock apoi lock pentru noi
        pthread_cond_wait(&thread_condition, &mutex);
    }
}

int main(int argc, char *argv[]) {
    pthread_t producer;

    // 250 msec
    delay.tv_sec  = 0;
    delay.tv_nsec = 250000000;

    pthread_mutex_init(&mutex, NULL);
    pthread_cond_init(&thread_condition, NULL);

    pthread_create(&producer, NULL, (void *)&producer_function, NULL);

    consumer_function();
}
\end{lstlisting}

\bigskip
\bigskip
\bigskip
\bigskip
\bigskip
\bigskip
\bigskip
\bigskip
\bigskip

\begin{figure}[H]
    \centering
    \includegraphics[width=0.8\textwidth]{3.png}
    \caption{Demonstratie \textbf{prodcons} inainte de modificari}
    \label{fig:3}
\end{figure}

\begin{figure}[H]
    \centering
    \includegraphics[width=0.8\textwidth]{4.png}
    \caption{Demonstratie \textbf{prodcons} dupa modificari}
    \label{fig:4}
\end{figure}

\newpage
\subsection*{Problema 4}
\begin{lstlisting}
#include <pthread.h>
#include <stdio.h>
#include <stdlib.h>

#define ITEMS 10

long buffer[ITEMS];
int  head = 0, tail = 0;

pthread_mutex_t mutex;
pthread_cond_t  thread_condition;
struct timespec delay;

// functia produce ia ca param `prod_id` care ne ajuta sa determinam din ce functie s-a apelat productia
long produce_item(int prod_id) {
    long item = random() % 256;
    printf("producing %d on id: %d\n", item, prod_id);

    return item;
}

void consume_item(long item) {
    printf("consuming %d\n", item);
}

void producer_function1(void) {
    while (1) {
        pthread_mutex_lock(&mutex);
        if ((tail + 1) % ITEMS != head) {
            buffer[tail] = produce_item(1);
            tail         = (tail + 1) % ITEMS;
        }
        pthread_mutex_unlock(&mutex);

        pthread_cond_broadcast(&thread_condition);
        nanosleep(&delay, NULL);
    }
}

void producer_function2(void) {
    while (1) {
        pthread_mutex_lock(&mutex);
        if ((tail + 1) % ITEMS != head) {
            buffer[tail] = produce_item(2);
            tail         = (tail + 1) % ITEMS;
        }
        pthread_mutex_unlock(&mutex);

        // facem broadcast pentru ca avem doi consumeri. refolosim, in mare, codul din exercitiul trecut
        pthread_cond_broadcast(&thread_condition);
        nanosleep(&delay, NULL);
    }
}

void consumer_function(void) {
    while (1) {
        if (head != tail) {
            consume_item(buffer[head]);
            head = (head + 1) % ITEMS;
        }
        // putem sa comentam lock/unlock pentru ca `cond_wait` o sa faca unlock apoi lock pentru noi
        pthread_cond_wait(&thread_condition, &mutex);
    }
}

int main(int argc, char *argv[]) {
    pthread_t producer;

    // 250 msec
    delay.tv_sec  = 0;
    delay.tv_nsec = 250000000;

    pthread_mutex_init(&mutex, NULL);
    pthread_cond_init(&thread_condition, NULL);

    // avem doua functii producator, 1 si 2, replici exacte ale functiei originale
    pthread_create(&producer, NULL, (void *)&producer_function1, NULL);
    pthread_create(&producer, NULL, (void *)&producer_function2, NULL);

    // functia consumer nu sufera modificari
    consumer_function();
    // cleanup
    pthread_cond_destroy(&thread_condition);
}

\end{lstlisting}


\begin{figure}[H]
    \centering
    \includegraphics[width=0.8\textwidth]{5.png}
    \caption{Demonstratie \textbf{prodcons} cu mai multi producatori}
    \label{fig:5}
\end{figure}


\newpage
\subsection*{Problema 5}
\begin{lstlisting}
#include <pthread.h>
#include <semaphore.h>
#include <stdio.h>
#include <stdlib.h>

#define ITEMS 10

long buffer[ITEMS];
int  head = 0, tail = 0;

sem_t free_slots, full_slots;

pthread_mutex_t mutex;
struct timespec delay;

long produce_item(void) {
    long item = random() % 256;
    printf("producing %d\n", item);

    return item;
}

// functia ia ca param `consumer_id` ca sa ne ajute sa identificam din ce consumer a fost chemata
void consume_item(long item, int consumer_id) {
    printf("consuming %d from id: %d\n", item, consumer_id);
}

void producer_function(void) {
    while (1) {
        sem_wait(&free_slots);
        pthread_mutex_lock(&mutex);

        if ((tail + 1) % ITEMS != head) {
            buffer[tail] = produce_item();
            tail         = (tail + 1) % ITEMS;
        }

        pthread_mutex_unlock(&mutex);
        sem_post(&full_slots);

        nanosleep(&delay, NULL);
    }
}

void consumer_function1(void) {
    while (1) {
        sem_wait(&full_slots);
        pthread_mutex_lock(&mutex);

        if (head != tail) {
            consume_item(buffer[head], 1);
            head = (head + 1) % ITEMS;
        }

        pthread_mutex_unlock(&mutex);
        sem_post(&free_slots);
    }
}

void consumer_function2(void) {
    while (1) {
        sem_wait(&full_slots);
        pthread_mutex_lock(&mutex);

        if (head != tail) {
            consume_item(buffer[head], 2);
            head = (head + 1) % ITEMS;
        }

        pthread_mutex_unlock(&mutex);
        sem_post(&free_slots);
    }
}

int main(int argc, char *argv[]) {
    pthread_t producer;
    // ave, nevoie de un thread pentru un consumator aditional
    pthread_t consumer_thread;

    // 250 msec
    delay.tv_sec  = 0;
    delay.tv_nsec = 250000000;

    sem_init(&free_slots, 0, ITEMS - 1);
    sem_init(&full_slots, 0, 0);

    pthread_mutex_init(&mutex, NULL);
    pthread_create(&producer, NULL, (void *)&producer_function, NULL);

    // thread doar pentru consumer1, consumer2 ramane in main thread
    pthread_create(&consumer_thread, NULL, (void *)&consumer_function1, NULL);
    consumer_function2();
}

\end{lstlisting}

\begin{figure}[H]
    \centering
    \includegraphics[width=0.8\textwidth]{6.png}
    \caption{Demonstratie \textbf{semprodcons} cu mai multi consumeri}
    \label{fig:6}
\end{figure}

\end{document}
